% Extensión cuántica de la teoría general de la medición.
% Se integra en el mismo capítulo fuente de 01_medida_ambivalencia.

\section{Conditioning of sections and quantum updating}
\label{sec:medicion-cuantica-genealogica-rev2}
\index{measurement!conditioning of sections}
\index{quantum instrument!genealogical realization}

The preceding general theory places measurement in the joint state of
system, reader, and apparatus. The quantum realization adds a precise correspondence
among three levels: restriction of the section space,
spectral projection, and the completely positive instrument. Joint evolution
belongs to coupling; publication of the outcome belongs to the
reader; updating belongs to conditioning on the read fiber.

\subsection{Measurement context \(\mathfrak M_a=(U_a,q_a,\mathcal C_a,m_0)\), output map, and conditioned fiber \(\Omega_{a,y}\)}

Let \((\Omega_S,\Sigma_S)\) and \((\Omega_M,\Sigma_M)\) be the measurable spaces
of genealogical states of the system and apparatus. A context
\(a\) is specified by
\begin{equation}
 \mathfrak M_a=(U_a,q_a,\mathcal C_a,m_0),
 \label{eq:contexto-medida-rev2}
\end{equation}
where \(m_0\in\Omega_M\) is the preparation of the apparatus,
\[
 U_a:\Omega_S\times\Omega_M
 \longrightarrow\Omega_S\times\Omega_M
\]
is a measurable coupling, \((Y_a,\Sigma_{Y_a})\) is the measurable space of
outcomes, \(q_a:\Omega_M\to Y_a\) is a measurable reader of the pointer, and
\(\mathcal C_a\) declares compatibility and the horizon of prolongation.
When a reversible regime is invoked, it will also be required that \(U_a\) be
bijective and bimeasurable; in the Hilbert realization, the
corresponding condition is the unitarity of \(\widehat U_a\).
The output is
\begin{equation}
 \operatorname{Out}_a(x)
 =q_a\!\left(\operatorname{pr}_M U_a(x,m_0)\right).
 \label{eq:salida-contexto-medida-rev2}
\end{equation}

Let \(\Omega_a\in\Sigma_S\) be the set of sections compatible with
the preparation and the context. The map \(\operatorname{Out}_a\) is
measurable by composition. For \(y\in Y_a\) with
\(\{y\}\in\Sigma_{Y_a}\), define the measurable fiber
\begin{equation}
 \Omega_{a,y}
 :=\{x\in\Omega_a:\operatorname{Out}_a(x)=y\}.
 \label{eq:fibra-resultado-medida-rev2}
\end{equation}

\begin{definition}[Genealogical Update]
The update determined by the result \(y\) is the restriction
\[
 \Omega_a\longmapsto\Omega_{a,y}.
\]
At finite depth it preserves exactly the prefixes that admit some
extension in \(\Omega_{a,y}\).
\end{definition}

For a prefix \(s\in S_n\), the compatible futures subsequent to the
reading form
\[
 \operatorname{Fut}_{a,y}(s)
 =\{x\in\Omega_{a,y}:x|_n=s\}.
\]
The update modifies the prolongation condition and preserves the already
traversed prefix together with the coupling record.

\begin{proposition}[Conditional probability over the fiber]
\label{prop:probabilidad-condicionada-medida-rev2}
Let \(\mu_a\) be a probability measure on the restricted measurable space
\((\Omega_a,\Sigma_S|_{\Omega_a})\), let
\(B\in\Sigma_S|_{\Omega_a}\), and suppose
\(\mu_a(\Omega_{a,y})>0\). Then
\begin{equation}
 \mu_{a,y}(B)
 =\frac{\mu_a(B\cap\Omega_{a,y})}
        {\mu_a(\Omega_{a,y})}
 \label{eq:bayes-genealogico-rev2}
\end{equation}
defines a probability concentrated on the fiber of the result:
\(\mu_{a,y}(\Omega_{a,y})=1\).
\end{proposition}

\begin{proof}
The map is nonnegative and countably additive by the properties of
\(\mu_a\). Normalization follows by substituting
\(B=\Omega_{a,y}\) into \eqref{eq:bayes-genealogico-rev2}.
\end{proof}

\subsection{Spectral realization of the device}

Let \(\{|r\rangle\}\) be an orthonormal basis adapted to the system observable,
let \(|m_0\rangle\) be the apparatus preparation, and let
\(\{|m_r\rangle\}\) be an orthonormal family of pointer states. The
unitary representation of the coupling satisfies
\begin{equation}
 \widehat U_a(|r\rangle\otimes|m_0\rangle)
 =|r\rangle\otimes|m_r\rangle.
 \label{eq:premedicion-espectral-rev2}
\end{equation}
For \(|\psi\rangle=\sum_r c_r|r\rangle\), linearity produces
\[
 \widehat U_a(|\psi\rangle\otimes|m_0\rangle)
 =\sum_r c_r|r\rangle\otimes|m_r\rangle.
\]
The coupling thus constructs a correlation between system and pointer. The
selective reading of the outcome \(r\) is represented, for a density state
with \(\operatorname{tr}(P_r\rho)>0\), by the
Luders~\cite{Lueders1950} rule,
\begin{equation}
 \rho\longmapsto
 \rho_r=\frac{P_r\rho P_r}{\operatorname{tr}(P_r\rho)}.
 \label{eq:luders-genealogico-rev2}
\end{equation}

\begin{theorem}[Correspondence between fibers and projections]
\label{thm:fibra-proyeccion-medida-rev2}
Suppose that each fiber \(\Omega_{a,y}\) is a finite union of disjoint
cylinders of a common depth \(n\), that the realization
\(\mathcal Q\) carries those cylinders to orthogonal subspaces, and that \(P_y\)
is the sum of their cylinder projectors. Let
\(\mu_a(\Omega_{a,y})>0\), and let \(\Psi_n(\mu_a)\) be the state constructed by
\eqref{eq:estado-amplitudes-medida-rev2} from that same measure. Then
\[
 \mu_a(\Omega_{a,y})=\|P_y\Psi_n(\mu_a)\|^2,
\]
and the restriction to \(\Omega_{a,y}\) is represented by the normalized
projection \(P_y\Psi_n(\mu_a)/\|P_y\Psi_n(\mu_a)\|\). In the density-matrix chart, if
\(\operatorname{tr}(\rho P_y)>0\), let \(\mathcal A_n\) be the algebra of cylindrical events of
depth \(n\), and suppose that \(B\mapsto P_B\) and \(P_y\) belong
to the same PVM and that
\begin{equation}
 \mu_a(B)=\operatorname{tr}(\rho P_B)
 \qquad\text{for every }B\in\mathcal A_n.
 \label{eq:equivalencia-medida-pvm-rev2}
\end{equation}
Then the same conditional probability is represented by
\eqref{eq:luders-genealogico-rev2}: for every \(B\in\mathcal A_n\),
\[
 \mu_{a,y}(B)=\operatorname{tr}(\rho_yP_B).
\]
\end{theorem}

\begin{proof}
The fiber is a disjoint union of cylinders, and \(P_y\) is the orthogonal sum of
their projectors. The \cref{thm:born-cilindrico-rev2} identifies the quadratic
norm of each component with its measure; additivity gives the equality.
After projection, the coefficients of each cylinder are divided by
\(\sqrt{\mu_a(\Omega_{a,y})}\); their squares are therefore the
conditional probabilities of \eqref{eq:bayes-genealogico-rev2}.
In the density chart, membership in the same PVM gives
\(P_BP_y=P_{B\cap\Omega_{a,y}}\). By cyclicity of the trace and by
\eqref{eq:equivalencia-medida-pvm-rev2},
\[
 \operatorname{tr}(\rho_yP_B)
 =\frac{\operatorname{tr}(\rho P_{B\cap\Omega_{a,y}})}
        {\operatorname{tr}(\rho P_y)}
 =\frac{\mu_a(B\cap\Omega_{a,y})}{\mu_a(\Omega_{a,y})}.
\]
\end{proof}

\subsection{Instruments induced by coupling}

For a refined pointer basis
\(\{|m_{y,\alpha}\rangle\}_{y,\alpha}\), the operators
\begin{equation}
 M_{y,\alpha}=\langle m_{y,\alpha}|\widehat U_a|m_0\rangle,
 \qquad
 \sum_{y,\alpha} M_{y,\alpha}^\dagger M_{y,\alpha}=I
 \label{eq:kraus-acoplamiento-genealogico-rev2}
\end{equation}
form a Kraus representation~\cite{Kraus1983}. The effect
\(E_y=\sum_\alpha M_{y,\alpha}^\dagger M_{y,\alpha}\) determines the
probability and, when \(p(y\mid a,\rho)>0\), the posterior state:
\begin{equation}
 p(y\mid a,\rho)=\operatorname{tr}(\rho E_y),
 \qquad
 \rho_{a,y}
 =\frac{\sum_\alpha M_{y,\alpha}\rho M_{y,\alpha}^\dagger}
 {p(y\mid a,\rho)}.
 \label{eq:instrumento-kraus-genealogico-rev2}
\end{equation}
The non-selective transformation is
\[
 \mathcal E_a(\rho)=
 \sum_{y,\alpha}M_{y,\alpha}\rho M_{y,\alpha}^\dagger.
\]

\begin{proposition}[Unitary coupling and reduced channel]
\label{prop:acoplamiento-canal-reducido-rev2}
The joint dynamics \(\widehat U_a\) and the preparation \(|m_0\rangle\)
induce a completely positive trace-preserving channel on the
system. The pointer decomposition induces the instrument
\(\mathcal I_{a,y}(\rho)=
\sum_\alpha M_{y,\alpha}\rho M_{y,\alpha}^\dagger\). Conversely, every completely positive channel
admits a dilation by an auxiliary space and an
isometry.
\end{proposition}

\begin{proof}
The completeness identity in
\eqref{eq:kraus-acoplamiento-genealogico-rev2} preserves the trace, and the Kraus
form proves complete positivity. The converse assertion is the
Stinespring dilation theorem~\cite{Stinespring1955}.
\end{proof}

The auxiliary space of the dilation can linearly realize the memory and
register transported by TPK once an embedding that intertwines
that memory with the auxiliary space has been declared. Stinespring guarantees the dilation of the
channel; it does not by itself identify the two objects. The Lueders rule is the fine
projective case \(M_{y,1}=P_y\); degenerate readers and nonideal
apparatus are described by means of families \(M_{y,\alpha}\).

\subsection{Internal Observer and Completeness of Reading}

A realized observer is a subsystem capable of preparing a context, coupling
to a boundary, preserving a distinguishable record, and using it as a
condition of prolongation. The sequence
\begin{equation}
 (x,m_0)
 \xrightarrow{\ U_a\ }(x',m')
 \xrightarrow{\ q_a\ }y
 \xrightarrow{\ \mathrm{cond}\ }\Omega_{a,y}
 \label{eq:secuencia-observador-rev2}
\end{equation}
places the observer's reading frame and memory within the extended state.

The Ambivalence principle completes this description. The areal and
volumetric readings produce quotients \(q_A,q_V\), whereas memory
\(T\) preserves transport between them. When
\[
 \Xi:x\longmapsto(q_A(x),q_V(x),T(x))
\]
is injective modulo the declared gauge group, the pair of observables and
the memory reconstruct the class of the complete state.

\begin{theorem}[Genealogical Decomposition of Measurement]
\label{thm:descomposicion-medicion-rev2}
Under the hypotheses of
\cref{thm:fibra-proyeccion-medida-rev2,prop:acoplamiento-canal-reducido-rev2}, the
transition between evolution and outcome decomposes, in the regime in which
\(U_a\) is bijective and bimeasurable or \(\widehat U_a\) is unitary, as follows:
\[
 \text{reversible coupling}
 \longrightarrow
 \text{pointer reading}
 \longrightarrow
 \text{section conditioning}.
\]
This decomposition possesses two regimes:
\begin{enumerate}
\item In a projective cylindrical measurement, \(M_{y,1}=P_y\) and the
fiber restriction, normalized projection, and Luders update coincide:
\[
 \Omega_a\longmapsto\Omega_{a,y},
 \qquad
 \Psi_n\longmapsto\frac{P_y\Psi_n}{\|P_y\Psi_n\|},
 \qquad
 \rho\longmapsto
 \frac{P_y\rho P_y}{\operatorname{tr}(\rho P_y)}.
\]
\item For a general instrument, the conditioned state is
\[
 \rho\longmapsto
 \frac{\sum_\alpha M_{y,\alpha}\rho M_{y,\alpha}^\dagger}
 {\operatorname{tr}(\rho E_y)}.
\]
In a dilation, this update arises from a pointer projection on the system
composed with the apparatus. Its correspondence with \(\Omega_{a,y}\) requires
a declared morphism intertwining the reader fiber, the effect
\(E_y\), and the instrument \(\mathcal I_{a,y}\).
\end{enumerate}
\end{theorem}

\begin{proof}
The first arrow realizes the coupling; the third restricts the space of
sections to the fiber defined in \eqref{eq:fibra-resultado-medida-rev2}. The
\cref{thm:fibra-proyeccion-medida-rev2} proves coincidence of the three
forms in the projective regime. In the general regime,
\eqref{eq:instrumento-kraus-genealogico-rev2} proves the instrument update, and
Stinespring supplies its dilated realization. The additional intertwining
morphism is precisely the datum transporting that
update back to the section space.
\end{proof}

\begin{statusbox}[title={Result and residence}]
In the declared model, the quantum update is the spectral realization of the restriction of compatible futures. Reversibility is determined by the bimeasurability of \(U_a\) or the unitarity of \(\widehat U_a\); information loss, in the fibers of the reader or in a subsequent erasure; and the conditioned state, in \(\Omega_{a,y}\). The thermodynamic cost of resetting and the zero Landauer cost of reversible evolution remain in the chapter on information and entropy.
\end{statusbox}
